\documentclass[11pt]{article}
\usepackage[paperwidth=11in,paperheight=8.5in,margin=0.25in]{geometry}
\usepackage[T1]{fontenc}
\usepackage[utf8]{inputenc}
\usepackage[english]{babel}
\usepackage{lmodern}
\usepackage{microtype}
\usepackage{amsmath,amssymb}
\usepackage{xcolor}
\usepackage{tikz}
\usetikzlibrary{calc}
\hyphenpenalty=8000
\exhyphenpenalty=8000
\emergencystretch=2em
\definecolor{QCBlue}{HTML}{1F4E79}
\definecolor{QCGreen}{HTML}{2E6B5F}
\definecolor{QCOrange}{HTML}{B65B2A}
\definecolor{QCPurple}{HTML}{4B3F72}
\definecolor{QCLight}{HTML}{F4F6F8}
\definecolor{QCBorder}{HTML}{D0D6DC}
\definecolor{QCText}{HTML}{111111}
\newlength{\Margin} \setlength{\Margin}{0.25in}
\newlength{\Gap} \setlength{\Gap}{0.1in}
\newlength{\TitleH} \setlength{\TitleH}{0.8in}
\newlength{\HeaderH} \setlength{\HeaderH}{0.38in}
\newlength{\Pad} \setlength{\Pad}{0.1in}
\newlength{\UsableW}
\newlength{\CellW}
\newlength{\CellH}
\setlength{\UsableW}{\dimexpr\paperwidth-2\Margin\relax}
\setlength{\CellW}{5.2in}
\setlength{\CellH}{3.5in}
\newcommand{\QuadBox}[4]{%
  \node[anchor=north west, draw=QCBorder, line width=0.6pt, rounded corners=10pt, fill=white, minimum width=\CellW, minimum height=\CellH, inner sep=0pt] (B) at #1 {};
  \fill[#2, rounded corners=10pt] ([xshift=0pt,yshift=0pt]B.north west) rectangle ([xshift=0pt,yshift=-\HeaderH]B.north east);
  \node[anchor=west, text=white, font=\bfseries] at ([xshift=0.18in,yshift=-0.22in]B.north west) {#3};
  \node[anchor=north west, text=QCText] at ([xshift=\Pad,yshift=-\HeaderH-\Pad]B.north west) {\begin{minipage}[t]{\dimexpr\CellW-2\Pad\relax}\raggedright\small #4\end{minipage}};
}
\pagestyle{empty}
\begin{document}
\begin{tikzpicture}[remember picture,overlay]
\coordinate (NW) at ([xshift=\Margin,yshift=-\Margin]current page.north west);
\node[anchor=north west, draw=QCBorder, line width=0.6pt, rounded corners=12pt, fill=QCLight, minimum width=\UsableW, minimum height=\TitleH, inner sep=0pt] (T) at (NW) {};
\node[anchor=north west] at ([xshift=0.22in,yshift=-0.22in]T.north west) {\begin{minipage}[t]{\dimexpr\UsableW-0.44in\relax}{\Huge\bfseries \textcolor{QCBlue}{Green's Functions Formula Sheet}}\par\vspace{2pt}{\normalsize \textbf{Diego Juarez}\hfill \textbf{Computational Electrodynamics}\hfill \textbf{\today}}\end{minipage}};
\coordinate (GridNW) at ([yshift=-0.15in]T.south west);
\coordinate (TL) at (GridNW);
\coordinate (TR) at ([xshift=\CellW+\Gap]GridNW);
\coordinate (BL) at ([yshift=-\CellH-\Gap]GridNW);
\coordinate (BR) at ([xshift=\CellW+\Gap,yshift=-\CellH-\Gap]GridNW);
\QuadBox{(TL)}{QCBlue}{1) Definitions}{
Free-space convention:
\[
\boxed{\nabla^2 G(\mathbf r,\mathbf r')=-4\pi\delta^{(3)}(\mathbf r-\mathbf r')}
\]
Field point: $\mathbf r$

Source point: $\mathbf r'$

Free-space Green function:
\[
G(\mathbf r,\mathbf r')=\frac{1}{|\mathbf r-\mathbf r'|}
\]
It is singular as $\mathbf r\to\mathbf r'$.
}
\QuadBox{(TR)}{QCGreen}{2) Representation Formula}{
Using Green's second identity:
\[
\Phi(\mathbf r)=\frac{1}{4\pi\varepsilon_0}\int_V \rho(\mathbf r')G(\mathbf r,\mathbf r')\,d^3r'
\]
\[
-\frac{1}{4\pi}\oint_S\left[
G\frac{\partial \Phi}{\partial n'}-\Phi\frac{\partial G}{\partial n'}
\right] da'
\]
Dirichlet Green function:
\[
G_D|_S=0
\]
Neumann Green function:
\[
\frac{\partial G_N}{\partial n'}\Big|_S=\text{specified}
\]
}
\QuadBox{(BL)}{QCOrange}{3) Symmetry and Examples}{
Reciprocity:
\[
G(\mathbf r,\mathbf r')=G(\mathbf r',\mathbf r)
\]
Grounded plane $z=0$:
\[
G_D=\frac{1}{|\mathbf r-\mathbf r'|}-\frac{1}{|\mathbf r-\mathbf r'^*|}
\]
where $\mathbf r'^*$ is the mirrored source.

Point charge near grounded plane:
\[
\Phi(\mathbf r)=\frac{q}{4\pi\varepsilon_0}\left(
\frac{1}{|\mathbf r-\mathbf r_0|}-\frac{1}{|\mathbf r-\mathbf r_0^*|}
\right)
\]
}
\QuadBox{(BR)}{QCPurple}{4) Practical Checks}{
\textbf{Always verify}

1. Correct delta-source equation.

2. Correct boundary condition on $S$.

3. Correct singularity near $\mathbf r=\mathbf r'$.

4. Reciprocity when applicable.

\textbf{Interpretation}

Green's function is the response kernel to a unit point source.

Different boundaries imply different Green's functions even for the same differential operator.
}
\end{tikzpicture}
\end{document}
